\section*{3. Theoretical Framework}

\subsection*{3.1 Generative AI as an Algorithmic Meaning-Making Agent}

This proposal treats generative AI as a special kind of information designer. Its difference from traditional recommendation algorithms lies not simply in whether it filters information, but in the fact that it generates a seemingly complete explanatory structure in continuous natural language. Traditional recommendation algorithms mainly alter the ordering and visibility of discrete information items. Generative AI, by contrast, compresses evidence selection, evidence weighting, conflict handling, and conclusion generation into a single answer, so that users confront an ``already interpreted'' information environment before they ever see the raw evidence.

More importantly, the conversational interface of generative AI induces users to make inferences about the system's knowledge state. Users do not merely see an answer; they may believe that the answer rests on extensive evidence processing, cross-dimensional comparison, and absorption of counterevidence. The informational effect of generative AI therefore lies not only in changing what users see, but also in changing what users believe the system has already seen, understood, and synthesized.

In traditional belief formation, the usual order is:
\[
\text{Evidence} \rightarrow \text{Explanation} \rightarrow \text{Judgment}.
\]
In a generative AI environment, the order is closer to:
\[
\text{AI Judgment} \rightarrow \text{AI Explanation} \rightarrow \text{Optional Evidence Verification}.
\]
This reversal in sequence is the key to the present theoretical framework. Users may encounter what looks like a complete posterior judgment first and only afterward decide whether to continue engaging with the underlying evidence. As a result, cognitive closure may occur before uncertainty has been fully confronted.

\subsection*{3.2 Algorithmic Authority and the Perception of Evidentiary Sufficiency}

This proposal defines \emph{algorithmic authority} as users' subjective perception that AI has relatively strong capability in information search, evidence integration, and judgment generation. Algorithmic authority is not the same as simple trust. Simple trust asks whether ``the AI's answer is correct.'' Algorithmic authority asks whether ``the AI has already processed the evidence environment sufficiently on my behalf.''

This distinction is crucial. Even if users do not fully believe the AI's final conclusion, they may still believe that the AI has already retrieved and integrated the main information. In that case, the influence of AI may not appear as direct adoption of the answer, but as reduced willingness to continue searching. In other words, algorithmic authority first affects users' judgments about evidentiary sufficiency.

In decision-making under uncertainty, whether individuals continue to search for information depends on their assessment of the value of additional information. If they believe continued search might substantially alter their judgment, continued sampling is reasonable. If they believe the existing information is already sufficient, stopping is also reasonable. The problem is that a synthetic AI output may cause users to overestimate the sufficiency of the current informational state and thereby reduce the subjective value of additional information too early.

Accordingly, the core function of algorithmic authority is to alter users' answers to the following questions:
\begin{itemize}[leftmargin=1.5em]
\item Do I still need more evidence?
\item Is the current answer already comprehensive enough?
\item Is continued search still worth it?
\end{itemize}

This means that the influence of AI is not simply to push users toward one specific answer. Rather, it changes their judgment about whether the evidence is already sufficient. It affects the stopping rule in the process of belief formation, not merely the final choice.

\subsection*{3.3 Path One: The Illusion of Algorithmic Exhaustiveness and Distortion in Information Acquisition}

The first mechanism path occurs at the information-acquisition stage. This proposal defines the \emph{illusion of algorithmic exhaustiveness} as users' belief that the answer generated by AI has already covered, filtered, and integrated a sufficiently broad evidence space, even though the AI may not actually have exhausted the relevant evidence.

This illusion does not necessarily arise from blind trust in AI. It may arise from the formal features of AI output itself. Generative AI answers are often clearly structured, linguistically fluent, complete in argument, and inclined to present multiple factors in a summarized way. Precisely this form of completeness makes it easy for users to feel that the main information has already been processed.

The illusion of algorithmic exhaustiveness lowers users' perceived value of additional information. Users may believe that continuing to search will bring only repetitive evidence or evidence with low marginal value, rather than evidence that could substantially change the current judgment. As a result, they become more likely to stop searching early, sample fewer pieces of evidence, and move into final decision-making more quickly.

This mechanism chain can be summarized as follows:
\[
\begin{aligned}
&\text{AI answer-first synthesis}
\rightarrow
\text{algorithmic authority}
\rightarrow
\text{illusion of algorithmic exhaustiveness} \\
&\rightarrow
\text{decline in perceived value of additional information}
\rightarrow
\text{premature search stopping} \\
&\rightarrow
\text{endogenous undersampling of evidence.}
\end{aligned}
\]

This path explains how AI may cause users to ``see less.''

It is important to emphasize that such behavior need not take the form of subjective laziness or overt irrationality. If users believe that the AI has already completed a sufficiently comprehensive integration of information, then reducing search may appear reasonable from their perspective. The real problem is that the ``formal completeness'' of AI output may be misread as ``evidentiary completeness.''

\subsection*{3.4 Path Two: Ambiguity Compression and Distortion in Information Processing}

The second mechanism path occurs at the information-processing stage. This proposal defines \emph{ambiguity compression} as the process by which AI converts heterogeneous, conflicting, and probabilistic evidence into a highly coherent explanatory narrative, thereby reducing users' perception of unresolved uncertainty and evidentiary conflict.

Complex problems usually admit multiple possible explanations. Evidence may be incomplete, internally conflicting, or point toward different conclusions under different conditions. Traditionally, such ambiguity reminds users to keep their judgment open and to revise it when new evidence appears. Generative AI, however, often tends to produce answers that are clear, readable, and coherent. In order to generate a ``helpful'' answer, AI may place conflicting evidence inside a unified narrative, explain counterexamples as boundary conditions or exceptions, and treat them as secondary factors.

For example, the AI may say:
\begin{quote}
Although some counterexamples exist, the overall trend still supports this judgment.
\end{quote}
\begin{quote}
Although the evidence is not entirely consistent, the more reasonable explanation is \ldots
\end{quote}
\begin{quote}
These differences mainly arise from contextual factors and do not alter the core conclusion.
\end{quote}

Such formulations are not necessarily wrong, but they change users' psychological perception of conflicting evidence. Evidence that might otherwise have strong diagnostic force may be reinterpreted as a local anomaly, background noise, or information already absorbed by the master narrative.

As a result, AI may influence not only how much evidence users see, but also how they interpret the evidence they do see. Even if users later encounter new disconfirming information, they may underestimate the extent to which that information challenges their prior judgment, producing insufficient belief updating.

This mechanism chain can be summarized as follows:
\[
\begin{aligned}
&\text{AI answer-first synthesis}
\rightarrow
\text{coherent explanatory narrative}
\rightarrow
\text{ambiguity compression} \\
&\rightarrow
\text{reduced perceived diagnosticity of disconfirming evidence}
\rightarrow
\text{insufficient belief updating} \\
&\rightarrow
\text{deviation from Bayesian updating.}
\end{aligned}
\]

This path explains how AI may cause users to ``update less.''

\subsection*{3.5 Double Bayesian Distortion}

The central theoretical concept of this proposal is \emph{double Bayesian distortion}. This concept refers to the possibility that generative AI creates simultaneous distortions at two stages of human belief formation: insufficient sampling at the information-acquisition stage and insufficient updating at the information-processing stage.

The first distortion is a distortion in information acquisition. In normative terms, a rational decision-maker should continue searching for evidence whenever the expected value of additional information exceeds the search cost. If the AI's synthetic answer causes users to believe mistakenly that the evidence is already sufficient, and they stop searching before acquiring enough information, a distortion arises on the information-acquisition side.

The second distortion is a distortion in information processing. In the normative Bayesian framework, after observing new evidence, individuals should revise posterior beliefs according to that evidence's diagnosticity. If later disconfirming evidence has high diagnostic value, but users under-update because they previously accepted the AI's coherent explanation, a distortion arises on the information-processing side.

The theoretical significance of double Bayesian distortion lies in the fact that it distinguishes between two error sources that are often conflated:
\begin{itemize}[leftmargin=1.5em]
\item Error type one: users did not acquire enough evidence.
\item Error type two: users acquired the evidence but did not give it enough weight.
\end{itemize}

The first is an information-acquisition problem; the second is an information-processing problem. Both can produce distorted posterior beliefs and lower decision quality, but they involve different mechanisms and therefore require different governance responses. If the problem is undersampling, the intervention should encourage continued search or reveal the incompleteness of the evidence space. If the problem is under-updating, the intervention should heighten the salience and diagnosticity of conflicting evidence.

The theoretical framework of the present proposal can therefore be summarized as follows:
\begin{align*}
\text{Answer-first synthesis by generative AI}
&\downarrow \\
\text{Algorithmic authority}
&\downarrow
\end{align*}

\begin{center}
\begin{tabular}{p{0.42\linewidth} p{0.42\linewidth}}
\toprule
\textbf{Information-acquisition path} & \textbf{Information-processing path} \\
\midrule
Illusion of algorithmic exhaustiveness & Ambiguity compression \\
$\rightarrow$ lower perceived value of additional information & $\rightarrow$ lower diagnosticity of disconfirming evidence \\
$\rightarrow$ premature stopping of sampling & $\rightarrow$ insufficient belief updating \\
\bottomrule
\end{tabular}
\end{center}

\[
\Downarrow
\]
\[
\text{Double Bayesian distortion}
\]
\[
\Downarrow
\]
\[
\text{Posterior belief distortion / decline in decision performance}.
\]

\section*{4. Research Questions}

The overarching research question of this proposal is:

\begin{quote}
How does answer-first synthesis by generative AI reshape the process of human belief formation under uncertainty by changing evidence sampling, ambiguity perception, and belief updating?
\end{quote}

More specifically, the proposal asks the following questions:

\textbf{RQ1: Does answer-first synthesis by generative AI change users' perceived signal structure?}

This question concerns how users understand the evidence-generation process behind an AI summary. The key issue is not whether users fully believe the AI's final conclusion, but whether they believe the AI has seen more evidence, covered more evidentiary dimensions, processed the main disconfirming evidence, and generated a synthesis that can represent the overall evidence environment.

\textbf{RQ2: Does perceived signal structure reduce the subjective value of continued search?}

This question concerns the psychological mechanism at the information-acquisition stage. If users believe that the AI summary has already adequately covered and integrated the relevant evidence, they may believe that continued search would bring only repetitive or low-marginal-value information and therefore reduce the perceived value of additional evidence.

\textbf{RQ3: Does generative AI produce undersampling of evidence through perceived signal structure?}

This question concerns the core behavioral outcome. If users in the AI answer-first condition voluntarily sample less evidence, stop searching earlier, and this effect is explained by perceived signal structure and perceived information value, then generative AI is not merely affecting final judgment; it is changing users' rules for information acquisition.

\section*{5. Research Hypotheses}

\subsection*{5.1 The Perceived Signal Structure Path}

\textbf{H1:} Compared with traditional search-style or evidence-list-style information presentation, AI answer-first increases inflation in users' perceived signal structure.

More specifically, users in the AI answer-first condition will believe that the summary is based on more evidence, broader evidentiary dimensions, and a more representative informational basis, and they will be more likely to believe that the AI has already processed the main disconfirming evidence.

\textbf{H2:} Inflation in perceived signal structure lowers the perceived informational value of continued search.

The more users believe that the AI summary has already sufficiently covered and represented the overall evidence pool, the more likely they are to regard the marginal value of drawing additional evidence as low.

\textbf{H3:} AI answer-first leads users to reduce the number of evidence samples they voluntarily draw.

Compared with evidence-list, search-style, or no-summary conditions, users in the AI answer-first condition will voluntarily sample fewer additional pieces of evidence and stop information search earlier.

\textbf{H4:} Perceived signal structure and perceived information value mediate the effect of AI answer-first on evidence undersampling.

AI answer-first does not simply make users ``too lazy to search.'' Rather, by changing users' subjective understanding of the coverage, representativeness, and degree of disconfirming-evidence processing behind the AI summary, it lowers the perceived necessity of continued search and thereby produces undersampling.

The mechanism chain is:
\[
\begin{aligned}
&\text{AI answer-first}
\rightarrow
\text{inflation in perceived signal structure} \\
&\rightarrow
\text{decline in perceived value of additional information}
\rightarrow
\text{decline in voluntary sample size.}
\end{aligned}
\]

\subsection*{5.2 The Signal-Structure Calibration Path}

\textbf{H5:} Signal-structure calibration prompts can weaken the effect of AI answer-first on undersampling.

If the AI summary explicitly discloses that its evidentiary coverage is limited, does not represent the full evidence set, or warns that disconfirming evidence may remain unprocessed, then inflation in users' perceived signal structure will be weaker, the perceived value of continued search will be higher, and voluntary sample size will also increase.

The mechanism chain is:
\[
\begin{aligned}
&\text{signal-structure calibration prompt}
\rightarrow
\text{reduced inflation in perceived signal structure} \\
&\rightarrow
\text{higher perceived value of additional information}
\rightarrow
\text{larger voluntary sample size.}
\end{aligned}
\]

\subsection*{5.3 Exclusion Tests}

\textbf{H6:} The effect of AI answer-first on undersampling cannot be explained entirely by general algorithmic trust or anchoring on the initial answer.

After controlling for users' trust in AI accuracy, their confidence in the initial conclusion, and the direction and strength of the initial answer, perceived signal structure should still significantly predict the perceived value of continued search and the number of evidence samples drawn voluntarily.

\section*{6. Core Theoretical Contributions}

This proposal seeks to contribute to the existing literature on AI and human judgment in three ways.

First, it reconceptualizes generative AI from ``algorithmic advice'' to an ``algorithmic meaning-making agent.'' Existing research often focuses on whether users adopt algorithmic output, whereas the present project emphasizes that the influence of generative AI arises at an earlier stage of problem understanding and evidence organization. Generative AI does not merely provide an answer. It organizes evidence in advance, compresses conflict, and generates an interpretive framework for the user.

Second, the proposal decomposes the influence of AI on human judgment into two stages: information acquisition and information processing. By introducing the mechanisms of the illusion of algorithmic exhaustiveness and ambiguity compression, it distinguishes between two different deviations: users ``seeing less'' and users ``updating less.'' This distinction helps avoid reducing the effect of AI to trust, laziness, anchoring, or automation bias alone.

Third, the proposal introduces the theoretical framework of double Bayesian distortion and embeds the behavioral effects of generative AI inside the normative frameworks of information acquisition, optimal stopping, and Bayesian updating. This framework helps explain not only how AI changes individual belief formation, but also how its potential welfare consequences and governance responses should be evaluated.

In short, the key influence of generative AI is not merely that it changes which answer users accept. It changes how users judge whether evidence is sufficient, whether conflict matters, and whether beliefs need revision. Generative AI is transforming humans from agents who form judgments on the basis of evidence into agents who manage evidence on the basis of algorithmic explanation. The present proposal studies the behavioral mechanisms and decision consequences of this shift in cognitive order.
